\allowdisplaybreaks
% \usepackage{cite}
\usepackage{graphicx}
\usepackage{tikz}
\usepackage{tikz-cd}
\usepackage{times}
\renewcommand\ttdefault{cmtt} 
\usepackage{bm}
\usepackage{physics}
\usepackage{xcolor}
\usepackage{natbib}
\usepackage{mdframed}
\usepackage{nicefrac}
\usepackage{booktabs}
\usepackage{lipsum}
\usepackage{titlesec}
\usepackage{wrapfig,lipsum,booktabs}
\usepackage{authblk}
\usepackage{blindtext}
\usepackage{subcaption}

\usepackage{algorithm}
\usepackage{algpseudocode}
% Algorithm definitions
\newcommand{\commentsymbol}{//}% or \% or $\triangleright$
\algrenewcommand\algorithmiccomment[1]{\hfill {\footnotesize \commentsymbol{} #1}}




\usepackage{titletoc}
\usepackage{todonotes}
\usepackage{authblk}
\usepackage{setspace}
\usepackage{dsfont} % for \mathbb{1} as \mathds{1}
\newcommand{\one}{\mathds{1}}

%%% Section Notation %%%
\usepackage[nameinlink,capitalize,noabbrev]{cleveref}
% \crefformat{section}{\S#2#1#3} % see manual of cleveref, section 8.2.1
% \crefformat{subsection}{\S#2#1#3}
% \crefformat{subsubsection}{\S#2#1#3}
% \crefformat{appendix}{\S #2#1#3}
% use: See \cref{sec:section}, \cref{sec:subsection} or \cref{sec:subsubsection}
%%% Section Notation %%%

%%% Chinese ###
\usepackage{CJK}
\newcommand{\zh}[1]{\begin{CJK}{UTF8}{bsmi}\normalfont\small #1\end{CJK}} % 繁體中文
\newcommand{\zhs}[1]{\begin{CJK}{UTF8}{gbsn}\normalfont \small#1\end{CJK}}% 簡體中文 
%%%% \begin{CJK}{UTF8}{bsmi}\normalfont 中文\end{CJK} when word missing %%%

%%%Claim%%%
\newenvironment{claim}{  \begin{mdframed}[linecolor=black!0,backgroundcolor=black!10]\noindent%\itshape
		\ignorespaces}{\end{mdframed}}

% \let\originalfigure=\figure
% \let\endoriginalfigure=\endfigure

% \renewenvironment{figure}[1][]{
%   \begin{originalfigure}[#1]
%     \begin{mdframed}[linecolor=black!0,backgroundcolor=black!1]
% }{
%     \end{mdframed}
%   \end{originalfigure}
% }
%%%Claim%%%

%%%% skip line in table %%%
\usepackage{array}
\usepackage{bbm}
\usepackage{makecell}

\renewcommand\theadalign{bc}
\renewcommand\theadfont{\bfseries}
\renewcommand\theadgape{\Gape[4pt]}
\renewcommand\cellgape{\Gape[0pt]}
%%%% skip line in table %%%

%% Comment%%
\newcommand{\jerry}[1]{\textcolor{red}{[JH: #1]}}
% \newcommand{\matt}[1]{\textcolor{firebrick}{[Matthew: #1]}}
\newcommand{\smj}[1]{\textcolor{blue}{[MJ: #1]}}
\newcommand{\zqj}[1]{\textcolor{red}{QZ: #1}}
% \newcommand{\juihui}[1]{\textcolor{magenta}{[juihui: #1]}}



\def\half{{\frac{1}{2}}}

%

\usepackage{dashbox}
\usepackage{xcolor}
\usepackage{colortbl}
\definecolor{lightyellow}{rgb}{1.0, 0.95, 0.7}
% \definecolor{Blue}{rgb}{0, 0.18, 0.65}
\definecolor{Blue}{rgb}{0, 0, 0.8}
\definecolor{blue}{rgb}{0,0,1}
% \definecolor{darkgreen}{rgb}{0.01,0.50,0.2}
\definecolor{darkgreen}{rgb}{0,0.40,0}
\definecolor{firebrick}{rgb}{0.698,0.133,0.133}
\newcommand*{\red}[1]{\textcolor{red}{#1}}
\newcommand*{\Red}[1]{\textcolor{firebrick}{#1}}
\newcommand*{\Blue}[1]{\textcolor{Blue}{#1}}
\newcommand*{\blue}[1]{\textcolor{blue}{#1}}
\newcommand*{\green}[1]{\textcolor{darkgreen}{#1}}
\newcommand*{\black}[1]{\textcolor{black}{#1}}

\definecolor{colorA}{rgb}{1,0,0}
\definecolor{colorB}{rgb}{0,0.3,1}
\definecolor{colorC}{rgb}{0.9,0.8,0.2}
\definecolor{colorD}{rgb}{0,0.65,0}
\definecolor{lesslightgray}{rgb}{0.5,0.5,0.5}
\definecolor{light-gray}{gray}{0.95}
\newcommand{\code}[1]{\colorbox{light-gray}{\footnotesize\texttt{#1}}}
\newcommand*{\grey}[1]{\textcolor{black!50}{#1}}
\newcommand*{\colorA}[1]{\textcolor{colorA}{#1}}
\newcommand*{\colorB}[1]{\textcolor{colorB}{#1}}
\newcommand*{\colorC}[1]{\textcolor{colorC}{#1}}
\newcommand*{\colorD}[1]{\textcolor{colorD}{#1}}


\let\tilde\widetilde
\let\hat\widehat

\newcommand{\calA}{\mathcal{A}}
\newcommand{\calB}{\mathcal{B}}
\newcommand{\calC}{\mathcal{C}}
\newcommand{\calD}{\mathcal{D}}
\newcommand{\calE}{\mathcal{E}}
\newcommand{\calF}{\mathcal{F}}
\newcommand{\calG}{\mathcal{G}}
\newcommand{\calH}{\mathcal{H}}
\newcommand{\calI}{\mathcal{I}}
\newcommand{\calJ}{\mathcal{J}}
\newcommand{\calK}{\mathcal{K}}
\newcommand{\calL}{\mathcal{L}}
\newcommand{\calM}{\mathcal{M}}
\newcommand{\calN}{\mathcal{N}}
\newcommand{\calO}{\mathcal{O}}
\newcommand{\calP}{\mathcal{P}}
\newcommand{\calQ}{\mathcal{Q}}
\newcommand{\calR}{\mathcal{R}}
\newcommand{\calS}{\mathcal{S}}
\newcommand{\calT}{\mathcal{T}}
\newcommand{\calU}{\mathcal{U}}
\newcommand{\calV}{\mathcal{V}}
\newcommand{\calW}{\mathcal{W}}
\newcommand{\calX}{\mathcal{X}}
\newcommand{\calY}{\mathcal{Y}}
\newcommand{\calZ}{\mathcal{Z}}




% \newcommand{\bA}{\mathbf{A}}
% \newcommand{\bB}{\mathbf{B}}
% \newcommand{\bC}{\mathbf{C}}
% \newcommand{\bD}{\mathbf{D}}
% \newcommand{\bE}{\mathbf{E}}
% \newcommand{\bF}{\mathbf{F}}
% \newcommand{\bG}{\mathbf{G}}
% \newcommand{\bH}{\mathbf{H}}
% \newcommand{\bI}{\mathbf{I}}
% \newcommand{\bJ}{\mathbf{J}}
% \newcommand{\bK}{\mathbf{K}}
% \newcommand{\bL}{\mathbf{L}}
% \newcommand{\bM}{\mathbf{M}}
% \newcommand{\bN}{\mathbf{N}}
% \newcommand{\bO}{\mathbf{O}}
% \newcommand{\bP}{\mathbf{P}}
% \newcommand{\bQ}{\mathbf{Q}}
% \newcommand{\bR}{\mathbf{R}}
% \newcommand{\bS}{\mathbf{S}}
% \newcommand{\bT}{\mathbf{T}}
% \newcommand{\bU}{\mathbf{U}}
% \newcommand{\bV}{\mathbf{V}}
% \newcommand{\bW}{\mathbf{W}}
% \newcommand{\bX}{\mathbf{X}}
% \newcommand{\bY}{\mathbf{Y}}
% \newcommand{\bZ}{\mathbf{Z}}
% \newcommand{\ba}{\mathbf{a}}
% \newcommand{\bb}{\mathbf{b}}
% \newcommand{\bc}{\mathbf{c}}
% \newcommand{\bd}{\mathbf{d}}
% \newcommand{\be}{\mathbf{e}}
% \newcommand{\bk}{\mathbf{k}}
% \newcommand{\bmf}{\mathbf{f}}
% \newcommand{\bu}{\mathbf{u}}
% \newcommand{\bp}{\mathbf{p}}
% \newcommand{\bs}{\mathbf{s}}
% \newcommand{\bg}{\mathbf{g}}
% \newcommand{\bh}{\mathbf{h}}
% \newcommand{\br}{\mathbf{r}}
% \newcommand{\bv}{\mathbf{v}}
% \newcommand{\bq}{\mathbf{q}}
% \newcommand{\bx}{\mathbf{x}}
% \newcommand{\by}{\mathbf{y}}
% \newcommand{\bz}{\mathbf{z}}
% \newcommand{\bxi}{{\bm{\xi}}}
\newcommand{\diag}{\mathop{\rm{diag}}}
\newcommand{\Max}{\mathop{\rm Max}}
\newcommand{\Min}{\mathop{\rm Min}}
\newcommand{\argmin}{\mathop{\mathrm{argmin}}}  
\newcommand{\argmax}{\mathop{\mathrm{argmax}}}  
\newcommand{\Var}{\mathop{\mathrm{Var}}}
\newcommand{\Sigmoid}{\mathop{\rm{Sigmoid}}}
\newcommand{\Softmax}{\mathop{\rm{Softmax}}}
\newcommand{\Sparsemax}{\mathop{\rm{Sparsemax}}}
\newcommand{\Sign}{\mathop{\rm{Sign}}}
\newcommand{\lse}{\mathop{\rm{lse}}}



\newcommand{\sumN}{\sum_{i=1}^N}
\newcommand{\sumK}{\sum_{k=1}^K}
\newcommand{\sumq}{\sum_{k=1}^q}
\newcommand{\summ}{\sum_{i=1}^m}
\newcommand{\sumM}{\sum_{\mu=1}^M}
\newcommand{\sumn}{\sum_{i=1}^n}
\newcommand{\sumd}{\sum_{i=1}^d}
\newcommand{\sumT}{\sum_{t=1}^T}

\def\P{\mathbb{P}}
\def\R{\mathbb{R}}
\def\B{\mathbb{B}}
%\def\E{\mathbb{E}}
\DeclareMathOperator*{\E}{{\mathbb{E}}} %%Zhao: I added this
\def\I{\mathcal{I}}
\def\p{\partial }
\let\cite\citep 

%%%%%%%%%%%%%%%%%%%%%%%%%%%%%%%%%%%%%%%%%%
\usepackage{amsthm}
\usepackage[many]{tcolorbox}
\setlength{\topsep}{0.8\baselineskip}  % Reduce space before and after theorems
\setlength{\partopsep}{\topsep}        % Consistent spacing for partopsep
% Adjust proof environment spacing manually
\usepackage{etoolbox}
% half‑baseline gap before every proof
\BeforeBeginEnvironment{proof}{\par\vspace{-1\baselineskip}}
% half‑baseline gap after every proof (and end the paragraph)
\AfterEndEnvironment  {proof}{\par\vspace{-1.5\baselineskip}}





% Adjusting spacing for all theorem-like environments
\newtheoremstyle{theoremstyle}
  {.5\baselineskip} % Space above
  {.5\baselineskip} % Space below
  {}                  % No font styling for body
  {}                  % Indentation
  {\bfseries}        % Heading font
  {.}                 % Punctuation
  {1em}               % Space after heading
  {}                  % Extra customization


\theoremstyle{theoremstyle}
\newtheorem{theorem}{Theorem}[section]
\newtheorem{lemma}{Lemma}[section]
\newtheorem{corollary}{Corollary}[theorem]
\newtheorem{proposition}{Proposition}[section]
\newtheorem{definition}{Definition}[section]
\newtheorem{assumption}{Assumption}[section]
\newtheorem{remark}{Remark}[section]
\newtheorem{problem}{Problem}
\newtheorem{property}{Property}

\tcolorboxenvironment{theorem}{
  breakable,
  colback=black!10,
  colframe=white,%
  width=\dimexpr\linewidth+10pt\relax,%  
  enlarge left by=-5pt,%                  
  enlarge right by=-5pt,%                 
  boxsep=5pt,%            
  boxrule=0pt,
  left=0pt,right=0pt,top=0pt,bottom=0pt,
  sharp corners,
  before skip=0.5\baselineskip, % Reduced space above
  after skip=0.5\baselineskip,  % Reduced space below
  fonttitle=\bfseries, % Bold title for theorem
  coltitle=black % Black title color
}
\tcolorboxenvironment{remark}{
  blanker,
  breakable,
  before skip=.8\baselineskip,  % half a baseline above
  after  skip=.8\baselineskip   % half a baseline below
}

\tcolorboxenvironment{proposition}{
  breakable,
  colback=black!10,
  colframe=white,%
  width=\dimexpr\linewidth+10pt\relax,%  
  enlarge left by=-5pt,%                  
  enlarge right by=-5pt,%                 
  boxsep=5pt,%            
  boxrule=0pt,
  left=0pt,right=0pt,top=0pt,bottom=0pt,
  sharp corners,
  before skip=0.5\baselineskip, % Reduced space above
  after skip=0.5\baselineskip,  % Reduced space below
  fonttitle=\bfseries, % Bold title for theorem
  coltitle=black % Black title color
}

\tcolorboxenvironment{lemma}{
  breakable,
  colback=black!10,
  colframe=white,%
  width=\dimexpr\linewidth+10pt\relax,%  
  enlarge left by=-5pt,%                  
  enlarge right by=-5pt,%                 
  boxsep=5pt,%            
  boxrule=0pt,
  left=0pt,right=0pt,top=0pt,bottom=0pt,
  sharp corners,
  before skip=0.5\baselineskip, % Reduced space above
  after skip=0.5\baselineskip,  % Reduced space below
  fonttitle=\bfseries, % Bold title for theorem
  coltitle=black % Black title color
}

\tcolorboxenvironment{corollary}{
  breakable,
  colback=black!10,
  colframe=white,%
  width=\dimexpr\linewidth+10pt\relax,%  
  enlarge left by=-5pt,%                  
  enlarge right by=-5pt,%                 
  boxsep=5pt,%            
  boxrule=0pt,
  left=0pt,right=0pt,top=0pt,bottom=0pt,
  sharp corners,
  before skip=0.5\baselineskip, % Reduced space above
  after skip=0.5\baselineskip,  % Reduced space below
  fonttitle=\bfseries, % Bold title for theorem
  coltitle=black % Black title color
}

\tcolorboxenvironment{definition}{
  breakable,
  colback=black!10,
  colframe=white,%
  width=\dimexpr\linewidth+10pt\relax,%  
  enlarge left by=-5pt,%                  
  enlarge right by=-5pt,%                 
  boxsep=5pt,%            
  boxrule=0pt,
  left=0pt,right=0pt,top=0pt,bottom=0pt,
  sharp corners,
  before skip=0.5\baselineskip, % Reduced space above
  after skip=0.5\baselineskip,  % Reduced space below
  fonttitle=\bfseries, % Bold title for theorem
  coltitle=black % Black title color
}
\tcolorboxenvironment{assumption}{
  breakable,
  colback=black!10,
  colframe=white,%
  width=\dimexpr\linewidth+10pt\relax,%  
  enlarge left by=-5pt,%                  
  enlarge right by=-5pt,%                 
  boxsep=5pt,%            
  boxrule=0pt,
  left=0pt,right=0pt,top=0pt,bottom=0pt,
  sharp corners,
  before skip=0.5\baselineskip, % Reduced space above
  after skip=0.5\baselineskip,  % Reduced space below
  fonttitle=\bfseries, % Bold title for theorem
  coltitle=black % Black title color
}

\crefname{theorem}{Theorem}{Theorems}
\crefname{proposition}{Proposition}{Propositions}
\crefname{lemma}{Lemma}{Lemmas}
\crefname{corollary}{Corollary}{Corollaries}
\crefname{definition}{Definition}{Definitions}
\crefname{assumption}{Assumption}{Assumptions}
\crefname{remark}{Remark}{Remarks}
\crefname{problem}{Problem}{Problems}
\crefname{property}{Property}{property}
\newtheorem{hypothesis}{Hypothesis}
\tcolorboxenvironment{hypothesis}{
  breakable,
  colback=black!10,
  colframe=white,%
  width=\dimexpr\linewidth+10pt\relax,%  
  enlarge left by=-5pt,%                  
  enlarge right by=-5pt,%                 
  boxsep=5pt,%            
  boxrule=0pt,
  left=0pt,right=0pt,top=0pt,bottom=0pt,
  sharp corners,
  before skip=0.5\baselineskip, % Reduced space above
  after skip=0.5\baselineskip,  % Reduced space below
  fonttitle=\bfseries, % Bold title for theorem
  coltitle=black % Black title color
}
\crefname{hypothesis}{Hypothesis}{Hypothesises}

\newtheorem{fact}{Fact}
\tcolorboxenvironment{fact}{
  breakable,
  colback=black!10,
  colframe=white,%
  width=\dimexpr\linewidth+10pt\relax,%  
  enlarge left by=-5pt,%                  
  enlarge right by=-5pt,%                 
  boxsep=5pt,%            
  boxrule=0pt,
  left=0pt,right=0pt,top=0pt,bottom=0pt,
  sharp corners,
  before skip=0.5\baselineskip, % Reduced space above
  after skip=0.5\baselineskip,  % Reduced space below
  fonttitle=\bfseries, % Bold title for theorem
  coltitle=black % Black title color
}
\crefname{fact}{Fact}{Facts}

\newtheorem{example}{Example}
\tcolorboxenvironment{example}{
  breakable,
  colback=black!10,
  colframe=white,%
  width=\dimexpr\linewidth+10pt\relax,%  
  enlarge left by=-5pt,%                  
  enlarge right by=-5pt,%                 
  boxsep=5pt,%            
  boxrule=0pt,
  left=0pt,right=0pt,top=0pt,bottom=0pt,
  sharp corners,
  before skip=0.5\baselineskip, % Reduced space above
  after skip=0.5\baselineskip,  % Reduced space below
  fonttitle=\bfseries, % Bold title for theorem
  coltitle=black % Black title color
}
\crefname{example}{Example}{Examples}

\newtheorem{question}{Question}
\tcolorboxenvironment{question}{
  breakable,
  colback=black!10,
  colframe=white,%
  width=\dimexpr\linewidth+10pt\relax,%  
  enlarge left by=-5pt,%                  
  enlarge right by=-5pt,%                 
  boxsep=5pt,%            
  boxrule=0pt,
  left=0pt,right=0pt,top=0pt,bottom=0pt,
  sharp corners,
  before skip=0.5\baselineskip, % Reduced space above
  after skip=0.5\baselineskip,  % Reduced space below
  fonttitle=\bfseries, % Bold title for theorem
  coltitle=black % Black title color
}






\crefname{question}{Question}{Questions}
\numberwithin{equation}{section}
\numberwithin{theorem}{section}
\numberwithin{proposition}{section}
\numberwithin{definition}{section}
\numberwithin{lemma}{section}
\numberwithin{assumption}{section}
\numberwithin{remark}{section}

%=== footnote without marker ===
% \blfootnote[]{}
\usepackage{lipsum}

\newcommand\blfootnote[1]{%
  \begingroup
  \renewcommand\thefootnote{}\footnote{#1}%
  \addtocounter{footnote}{-1}%
  \endgroup
}
%=== footnote without marker ===
%%% derivation step description
\newcommand*{\annot}[1]{\tag*{\footnotesize{\textcolor{black!50}{\big(#1\big)}}}}

%%% Better Bar Notation%%%
\makeatletter
\let\save@mathaccent\mathaccent
\newcommand*\if@single[3]{%
    \setbox0\hbox{${\mathaccent"0362{#1}}^H$}%
    \setbox2\hbox{${\mathaccent"0362{\kern0pt#1}}^H$}%
    \ifdim\ht0=\ht2 #3\else #2\fi
}
%The bar will be moved to the right by a half of \macc@kerna, which is computed by amsmath:
\newcommand*\rel@kern[1]{\kern#1\dimexpr\macc@kerna}
%If there's a superscript following the bar, then no negative kern may follow the bar;
%an additional {} makes sure that the superscript is high enough in this case:
\newcommand*\widebar[1]{\@ifnextchar^{{\wide@bar{#1}{0}}}{\wide@bar{#1}{1}}}
%Use a separate algorithm for single symbols:
\newcommand*\wide@bar[2]{\if@single{#1}{\wide@bar@{#1}{#2}{1}}{\wide@bar@{#1}{#2}{2}}}
\newcommand*\wide@bar@[3]{%
    \begingroup
    \def\mathaccent##1##2{%
        %Enable nesting of accents:
        \let\mathaccent\save@mathaccent
        %If there's more than a single symbol, use the first character instead (see below):
        \if#32 \let\macc@nucleus\first@char \fi
        %Determine the italic correction:
        \setbox\z@\hbox{$\macc@style{\macc@nucleus}_{}$}%
        \setbox\tw@\hbox{$\macc@style{\macc@nucleus}{}_{}$}%
        \dimen@\wd\tw@
        \advance\dimen@-\wd\z@
        %Now \dimen@ is the italic correction of the symbol.
        \divide\dimen@ 3
        \@tempdima\wd\tw@
        \advance\@tempdima-\scriptspace
        %Now \@tempdima is the width of the symbol.
        \divide\@tempdima 10
        \advance\dimen@-\@tempdima
        %Now \dimen@ = (italic correction / 3) - (Breite / 10)
        \ifdim\dimen@>\z@ \dimen@0pt\fi
        %The bar will be shortened in the case \dimen@<0 !
        \rel@kern{0.6}\kern-\dimen@
        \if#31
        \overline{\rel@kern{-0.6}\kern\dimen@\macc@nucleus\rel@kern{0.4}\kern\dimen@}%
        \advance\dimen@0.4\dimexpr\macc@kerna
        %Place the combined final kern (-\dimen@) if it is >0 or if a superscript follows:
        \let\final@kern#2%
        \ifdim\dimen@<\z@ \let\final@kern1\fi
        \if\final@kern1 \kern-\dimen@\fi
        \else
        \overline{\rel@kern{-0.6}\kern\dimen@#1}%
        \fi
    }%
    \macc@depth\@ne
    \let\math@bgroup\@empty \let\math@egroup\macc@set@skewchar
    \mathsurround\z@ \frozen@everymath{\mathgroup\macc@group\relax}%
    \macc@set@skewchar\relax
    \let\mathaccentV\macc@nested@a
    %The following initialises \macc@kerna and calls \mathaccent:
    \if#31
    \macc@nested@a\relax111{#1}%
    \else
    %If the argument consists of more than one symbol, and if the first token is
    %a letter, use that letter for the computations:
    \def\gobble@till@marker##1\endmarker{}%
    \futurelet\first@char\gobble@till@marker#1\endmarker
    \ifcat\noexpand\first@char A\else
    \def\first@char{}%
    \fi
    \macc@nested@a\relax111{\first@char}%
    \fi
    \endgroup
    }
\makeatother
\let\bar\widebar
